\section[Cote 161-1 : la partie fixe d'une adjonction, les adjonctions idempotentes, les objets locaux]{Cote 161-1 : la partie fixe d'une adjonction, les adjonctions idempotentes, les objets locaux\protect\verifie{Folder161_1}}\label{sec:161-1}

La cote 161-1, \emph{Catégories : notes manuscrites (s.d.)}, dix-neuf pages sans date, appartient au groupe des dossiers rassemblés par Grothendieck sur différentes thématiques, que l'inventaire date \q{[après 1961-vers 1977]}. Les pages~3 à~7 regardent ce qu'un couple de foncteurs adjoints $E \rightleftarrows F$ induit entre les deux catégories. La page~3 s'ouvre sur \q{Foncteurs adjoints [demander Schanuel]}, écrit le critère~(1) de pleine fidélité sur une sous-catégorie pleine et porte en marge \q{donner un exemple montrant que cette condition n'est pas surabondante, avec $E'$ réduit à un élément} (\lot{161-1}{1}{3}) ; la page~5 dresse la liste des conditions équivalentes qui définissent aujourd'hui une adjonction idempotente (\lot{161-1}{1}{5}) ; la page~7 se demande si une sous-catégorie réflexive se reconnaît comme lieu des objets locaux pour les flèches que le réflecteur inverse, et conclut \q{Or c'est vrai sauf erreur…} (\lot{161-1}{1}{7}). On note, comme la lecture, $u$ l'adjoint à gauche de $v$, $\eta : \mathrm{id}_E \to vu$ l'unité et $\varepsilon : uv \to \mathrm{id}_F$ la coünité (le manuscrit écrit $i$ et $j$).

\begin{enonce}[(1), page 3]
Soient $E' \subset E$ une sous-catégorie pleine et $\overline{E}'$ sa clôture par isomorphie. Alors
\[
  \begin{gathered}
  u|E' \text{ est pleinement fidèle et } vu(E') \subset \overline{E}' \\
  \iff \text{pour tout } Y \in E',\ \eta_Y : Y \to vu(Y) \text{ est un isomorphisme.}
  \end{gathered}
\]
La seconde condition du membre de gauche n'est pas surabondante : la pleine fidélité de $u|E'$ ne force pas $vu$ à ramener $E'$ dans $\overline{E}'$.
\end{enonce}

\begin{enonce}[Adjonctions idempotentes, page 5]
Les conditions \q{$\eta v : v \to vuv$ est un isomorphisme} et \q{$\varepsilon u : uvu \to u$ est un isomorphisme} de la condition~2 du manuscrit sont redondantes : chacune des quatre conditions ($\eta v$, $v\varepsilon$, $\varepsilon u$, $u\eta$ inversible) implique les trois autres.
\end{enonce}

\begin{enonce}[Objets locaux, pages 6 et 7]
Soient $\alpha : E_1 \to E$ un foncteur pleinement fidèle admettant un adjoint à gauche $\alpha'$, et $\Sigma$ la classe des flèches de $E$ que $\alpha'$ rend inversibles. Une flèche $f : X \to Y$ est dans $\Sigma$ si et seulement si $\operatorname{Hom}_E(Y, \alpha Z_1) \to \operatorname{Hom}_E(X, \alpha Z_1)$ est bijective pour tout $Z_1 \in E_1$ ; et un objet $\xi$ de $E$ tel que $\operatorname{Hom}_E(-, \xi)$ transforme les flèches de $\Sigma$ en bijections (un objet $\Sigma$-local) est dans l'image essentielle de $\alpha$.
\end{enonce}

On rappelle que la bijection d'adjonction $\operatorname{Hom}_F(uX, P) \simeq \operatorname{Hom}_E(X, vP)$ envoie $g$ sur $\eta_X \circ v(g)$ (on compose de gauche à droite : $\eta_X$, puis $v(g)$), et qu'on a les identités triangulaires $\eta_{vP}\,;v(\varepsilon_P) = \mathrm{id}_{vP}$ et $u(\eta_X)\,;\varepsilon_{uX} = \mathrm{id}_{uX}$. Dans les démonstrations qui suivent, on écrit $f\,;g$ pour « $f$ puis $g$ ».

\begin{lemme}\label{lem:161-1-hom}
Pour $X, Y \in E$, l'application $\operatorname{Hom}_E(X, Y) \to \operatorname{Hom}_F(uX, uY)$, $f \mapsto u(f)$, est bijective si et seulement si l'application $f \mapsto f\,;\eta_Y$, de $\operatorname{Hom}_E(X, Y)$ dans $\operatorname{Hom}_E(X, vuY)$, l'est.
\end{lemme}

\begin{proof}
La seconde application est la composée de la première et de la bijection d'adjonction : celle-ci envoie $u(f)$ sur $\eta_X\,;vu(f)$, qui vaut $f\,;\eta_Y$ par naturalité de $\eta$. C'est la phrase de la page~3 : l'application \q{se déduit de l'application $i(Y) : Y \to vu(Y)$ au moyen de $\mathrm{Hom}(X, i(Y))$}.
\end{proof}

\begin{proof}[Démonstration de (1)]
Supposons que $\eta_Y$ soit un isomorphisme pour tout $Y \in E'$. Pour $X, Y \in E'$, la composition avec $\eta_Y$ est alors bijective, et $u|E'$ est pleinement fidèle par le lemme~\ref{lem:161-1-hom} ; et $vu(Y)$, isomorphe à $Y$, est dans $\overline{E}'$.

Réciproquement, soit $Y \in E'$ ; par hypothèse il existe $Z \in E'$ et un isomorphisme $e : vu(Y) \to Z$. Le lemme~\ref{lem:161-1-hom}, appliqué aux couples $(Z, Y)$ et $(Y, Y)$ d'objets de $E'$, dit que $f \mapsto f\,;\eta_Y$ est bijective de $\operatorname{Hom}(Z, Y)$ sur $\operatorname{Hom}(Z, vuY)$ et de $\operatorname{Hom}(Y, Y)$ sur $\operatorname{Hom}(Y, vuY)$. Par la surjectivité de la première, il existe $g : Z \to Y$ avec $g\,;\eta_Y = e^{-1}$. Posons $s = e\,;g : vu(Y) \to Y$. D'une part $s\,;\eta_Y = e\,;e^{-1} = \mathrm{id}_{vuY}$. D'autre part $(\eta_Y\,;s)\,;\eta_Y = \eta_Y\,;(s\,;\eta_Y) = \eta_Y = \mathrm{id}_Y\,;\eta_Y$, et l'injectivité de la seconde bijection donne $\eta_Y\,;s = \mathrm{id}_Y$. Donc $\eta_Y$ est un isomorphisme, d'inverse $s$.
\end{proof}

\begin{proof}[L'exemple demandé en marge]
Prenons pour $E$ la catégorie des ensembles, pour $F$ la catégorie ponctuelle, pour $u$ l'unique foncteur $E \to F$ et pour $v$ le foncteur qui pointe un ensemble à un élément $\{*\}$ ; on a $u \dashv v$ parce que $\{*\}$ est final. Soit $E' = \{\varnothing\}$, réduite à un objet. Alors $u|E'$ est pleinement fidèle, puisque $\operatorname{Hom}(\varnothing, \varnothing)$ et $\operatorname{Hom}(u\varnothing, u\varnothing)$ ont chacun un seul élément ; mais $vu(\varnothing) = \{*\}$ n'est pas isomorphe à $\varnothing$, et $\eta_\varnothing : \varnothing \to \{*\}$ n'est pas un isomorphisme. La condition $vu(E') \subset \overline{E}'$ est donc nécessaire à l'équivalence~(1).
\end{proof}

\begin{proof}[Démonstration de l'énoncé sur les adjonctions idempotentes]
On note (a) $\eta_{vP}$ inversible pour tout $P$, (b) $v(\varepsilon_P)$ inversible pour tout $P$, (c) $\varepsilon_{uX}$ inversible pour tout $X$, (d) $u(\eta_X)$ inversible pour tout $X$. On utilise que si $f\,;g$ est une identité et que l'une de $f$, $g$ est un isomorphisme, l'autre en est un, inverse de la première.

(a)~$\Leftrightarrow$~(b) et (c)~$\Leftrightarrow$~(d) résultent de là et des identités triangulaires $\eta_{vP}\,;v(\varepsilon_P) = \mathrm{id}$ et $u(\eta_X)\,;\varepsilon_{uX} = \mathrm{id}$.

(a)~$\Rightarrow$~(c). Soit $X \in E$. Par (a)~$\Rightarrow$~(b) en $P = uX$, $v(\varepsilon_{uX})$ est inversible. Or $vu(\eta_X)$ et $\eta_{vuX}$ en sont deux sections : $vu(\eta_X)\,;v(\varepsilon_{uX}) = v\bigl(u(\eta_X)\,;\varepsilon_{uX}\bigr) = \mathrm{id}$, et $\eta_{vuX}\,;v(\varepsilon_{uX}) = \mathrm{id}$ par l'identité triangulaire ; elles sont donc égales. Montrons que $u(\eta_X)$ est inverse de $\varepsilon_{uX}$. On a déjà $u(\eta_X)\,;\varepsilon_{uX} = \mathrm{id}$. La naturalité de $\varepsilon$ appliquée à $u(\eta_X) : uX \to uvuX$ donne
\[
  \varepsilon_{uX}\,;u(\eta_X) = uv\bigl(u(\eta_X)\bigr)\,;\varepsilon_{uvuX} = u(\eta_{vuX})\,;\varepsilon_{uvuX} = \mathrm{id},
\]
la deuxième égalité par ce qui précède, la troisième par l'identité triangulaire en $vuX$.

(c)~$\Rightarrow$~(a) est l'énoncé dual. Soit $P \in F$. Par (c)~$\Rightarrow$~(d) en $X = vP$, $u(\eta_{vP})$ est inversible ; $uv(\varepsilon_P)$ et $\varepsilon_{uvP}$ en sont deux rétractions ($u(\eta_{vP})\,;uv(\varepsilon_P) = u(\eta_{vP}\,;v(\varepsilon_P)) = \mathrm{id}$, et l'identité triangulaire), donc sont égales. Alors $\eta_{vP}\,;v(\varepsilon_P) = \mathrm{id}$, et la naturalité de $\eta$ appliquée à $v(\varepsilon_P)$ donne $v(\varepsilon_P)\,;\eta_{vP} = \eta_{vuvP}\,;vuv(\varepsilon_P) = \eta_{vuvP}\,;v(\varepsilon_{uvP}) = \mathrm{id}$.

Les quatre conditions sont donc équivalentes ; en particulier la condition~2 du manuscrit, $\eta v$ \emph{et} $\varepsilon u$ inversibles, équivaut à sa première moitié.
\end{proof}

\begin{proof}[Démonstration de l'énoncé sur les objets locaux]
\emph{Caractérisation de $\Sigma$.} C'est un fait classique de la théorie des sous-catégories réflexives. Par l'adjonction $\alpha' \dashv \alpha$, $\operatorname{Hom}_E(Y, \alpha Z_1) \simeq \operatorname{Hom}_{E_1}(\alpha' Y, Z_1)$ naturellement en $Y$, de sorte que la composition avec $f$ à la source correspond à la composition avec $\alpha'(f)$. Par le lemme de Yoneda, $\alpha'(f)$ est un isomorphisme si et seulement si $\operatorname{Hom}_{E_1}(\alpha' Y, Z_1) \to \operatorname{Hom}_{E_1}(\alpha' X, Z_1)$ est bijective pour tout $Z_1$. En particulier tout objet $\alpha Z_1$, et tout objet qui lui est isomorphe, est $\Sigma$-local.

\emph{Un objet local est dans l'image essentielle.} Soit $\xi$ un objet $\Sigma$-local et $\eta = \eta_\xi : \xi \to \alpha\alpha'\xi$ l'unité. Pour tout $Z_1$, la composition avec $\eta$, de $\operatorname{Hom}_E(\alpha\alpha'\xi, \alpha Z_1)$ dans $\operatorname{Hom}_E(\xi, \alpha Z_1)$, est bijective : c'est la bijection d'adjonction $\operatorname{Hom}_{E_1}(\alpha'\xi, Z_1) \simeq \operatorname{Hom}_E(\xi, \alpha Z_1)$ lue à travers la pleine fidélité de $\alpha$. Par la caractérisation qui précède, $\eta \in \Sigma$. Comme $\xi$ est local, la composition avec $\eta$, de $\operatorname{Hom}(\alpha\alpha'\xi, \xi)$ dans $\operatorname{Hom}(\xi, \xi)$, est bijective ; il existe donc $r : \alpha\alpha'\xi \to \xi$ avec $\eta\,;r = \mathrm{id}_\xi$. Alors $\eta\,;(r\,;\eta) = \eta = \eta\,;\mathrm{id}$, et comme la composition avec $\eta$ est injective sur les flèches à valeurs dans l'objet local $\alpha\alpha'\xi$, $r\,;\eta = \mathrm{id}$. Ainsi $\xi \simeq \alpha(\alpha'\xi)$. Avec la dernière phrase du paragraphe précédent, on obtient que les objets $\Sigma$-locaux sont exactement ceux de l'image essentielle de $\alpha$.
\end{proof}

\begin{enonce}[(2), page 3]
Soient $F' \subset F$ une sous-catégorie pleine et $\overline{F}'$ sa clôture par isomorphie. Alors
\[
  \begin{gathered}
  v|F' \text{ est pleinement fidèle et } uv(F') \subset \overline{F}' \\
  \iff \text{pour tout } P \in F',\ \varepsilon_P : uv(P) \to P \text{ est un isomorphisme.}
  \end{gathered}
\]
\end{enonce}

\begin{proof}
C'est la démonstration de~(1) retournée. Pour $P, Q \in F$, l'application $g \mapsto v(g)$, de $\operatorname{Hom}_F(P, Q)$ dans $\operatorname{Hom}_E(vP, vQ)$, suivie de la bijection d'adjonction réciproque $h \mapsto u(h)\,;\varepsilon_Q$, donne $g \mapsto uv(g)\,;\varepsilon_Q = \varepsilon_P\,;g$ par naturalité de $\varepsilon$ : c'est le \q{$\mathrm{Hom}(j(P), Q)$} de la page, et $v$ est bijective sur $\operatorname{Hom}(P, Q)$ si et seulement si $g \mapsto \varepsilon_P\,;g$ l'est. Si $\varepsilon_P$ est inversible pour tout $P \in F'$, la composition avec $\varepsilon_P$ est bijective, et $uv(P) \simeq P$. Réciproquement, soient $P \in F'$, $Z \in F'$ et $e : uv(P) \to Z$ un isomorphisme ; la surjectivité de $g \mapsto \varepsilon_P\,;g$ sur $\operatorname{Hom}(P, Z)$ donne $g$ avec $\varepsilon_P\,;g = e$. Posons $s = g\,;e^{-1}$. Alors $\varepsilon_P\,;s = \mathrm{id}$, et $\varepsilon_P\,;(s\,;\varepsilon_P) = \varepsilon_P = \varepsilon_P\,;\mathrm{id}$ ; l'injectivité sur $\operatorname{Hom}(P, P)$ donne $s\,;\varepsilon_P = \mathrm{id}$.
\end{proof}

\begin{enonce}[Les cinq conditions, page 3]
Soient $E' \subset E$ et $F' \subset F$ des sous-catégories strictement pleines ; on note $u' : E' \to F'$ et $v' : F' \to E'$ les foncteurs induits quand $u(E') \subset F'$ et $v(F') \subset E'$. Les conditions suivantes sont équivalentes :
\begin{enumerate}
\item $E'$ vérifie~(1), $F'$ vérifie~(2), et $F' = \overline{u(E')}$, $E' = \overline{v(F')}$ ;
\item $u(E') \subset F'$, $v(F') \subset E'$, et $u'$, $v'$ sont pleinement fidèles ;
\item $u(E') \subset F'$, $v(F') \subset E'$, et $u'$ est une équivalence ;
\item $u(E') \subset F'$, $v(F') \subset E'$, et $v'$ est une équivalence.
\end{enumerate}
De plus, si $E'$ strictement pleine vérifie~(1), $\alpha(E') = \overline{u(E')}$ vérifie~(2) et $\beta(\alpha(E')) = E'$ ; dualement pour $F'$.
\end{enonce}

\begin{proof}
On lit \q{$u'$ est une équivalence} comme \q{$u'$ est pleinement fidèle et tout objet de $F'$ est isomorphe à un $u(Y)$, $Y \in E'$}. On utilise deux faits : si $\eta_Y$ est inversible, $\varepsilon_{uY}$ l'est, puisque $u(\eta_Y)\,;\varepsilon_{uY} = \mathrm{id}$ ; et si $\varepsilon_P$ l'est, $\eta_{vP}$ l'est, par l'autre identité triangulaire. Les objets où $\eta$ (resp. $\varepsilon$) est inversible forment une sous-catégorie strictement pleine, par naturalité.

Si $u(E') \subset F'$, $v(F') \subset E'$, $\eta$ est inversible sur $E'$ et $\varepsilon$ sur $F'$, alors 1.\ vaut : un $P \in F'$ est isomorphe, par $\varepsilon_P$, à $u(vP)$ avec $vP \in E'$, et tout objet isomorphe à un $u(Y)$, $Y \in E'$, est dans $F'$, strictement pleine ; de même pour $E'$, par $\eta_Y^{-1}$. Réciproquement 1.\ donne $u(E') \subset F'$, $v(F') \subset E'$, et, par~(1) et~(2), la pleine fidélité de $u'$ et $v'$ et l'essentielle surjectivité de $u'$ et de $v'$ ; donc 1.\ entraîne 2., 3.\ et~4.

Sous 2., $vu(E') \subset v(F') \subset E'$, et (1) donne $\eta$ inversible sur $E'$ ; de même (2) donne $\varepsilon$ inversible sur $F'$ ; d'où 1. Sous 3., (1) donne de même $\eta$ inversible sur $E'$ ; un $P \in F'$ est isomorphe à un $u(Y)$, $Y \in E'$, où $\varepsilon$ est inversible par le premier fait ; donc $\varepsilon_P$ l'est, d'où 1. Le cas~4.\ est symétrique.

Pour la bijection : si $Y \in E'$, $Y \simeq vu(Y)$ par $\eta_Y$ et $uY \in \alpha(E')$ ; si $Y \simeq v(P)$ avec $P \simeq u(Y')$, $Y' \in E'$, alors $Y \simeq vu(Y') \simeq Y'$ et $Y \in E'$. Et $\varepsilon$ est inversible sur $\alpha(E')$ par le premier fait.
\end{proof}

\begin{enonce}[Les points fixes, page 4]
Soient $E_0 \subset E$ la sous-catégorie pleine des $Y$ où $\eta_Y$ est inversible et $F_0 \subset F$ celle des $P$ où $\varepsilon_P$ l'est. Ce sont les plus grandes sous-catégories vérifiant~(1) et~(2) ; on a $F_0 = \overline{u(E_0)}$, $E_0 = \overline{v(F_0)}$, et $u$, $v$ induisent des équivalences quasi-inverses $E_0 \simeq F_0$.
\end{enonce}

\begin{proof}
La maximalité est la forme de~(1) et~(2) démontrée plus haut : $E'$ vérifie~(1) si et seulement si $E' \subset E_0$. Par les deux faits de la démonstration précédente, $u(E_0) \subset F_0$ et $v(F_0) \subset E_0$ ; les égalités $F_0 = \overline{u(E_0)}$ et $E_0 = \overline{v(F_0)}$ en résultent comme dans le cas~1.\ ci-dessus. Les foncteurs induits $u_0 : E_0 \to F_0$ et $v_0 : F_0 \to E_0$ sont quasi-inverses, par les isomorphismes naturels $\eta_Y : Y \to v_0u_0(Y)$ et $\varepsilon_P : u_0v_0(P) \to P$.
\end{proof}

\begin{enonce}[Conditions 1), 3) et~4), page 5 ; la préservation des colimites]
La condition~1), $\operatorname{Im} v \subset E_0$ et $\operatorname{Im} u \subset F_0$, est l'idempotence, et chacune de ses moitiés suffit. Si $\operatorname{Im} u \subset F_0$, l'inclusion $E_0 \subset E$ admet pour adjoint à gauche $v_0 \circ u'$, $X \mapsto vu(X)$ ; si $\operatorname{Im} v \subset E_0$, l'inclusion $F_0 \subset F$ admet pour adjoint à droite $u_0 \circ v'$, $P \mapsto uv(P)$. Ce réflecteur $E \to E_0$ commute aux limites inductives et l'inclusion $E_0 \subset E$ aux limites projectives ; $E_0$ a les limites inductives qu'a $E$ ; dualement, le coréflecteur $F \to F_0$ commute aux limites projectives, l'inclusion $F_0 \subset F$ aux limites inductives, et $F_0$ a les limites projectives qu'a $F$. Enfin, pour une adjonction idempotente, $u \simeq \beta\, u_0\, \alpha'$ et $v \simeq \alpha\, v_0\, \beta'$, avec $\alpha'$ le réflecteur, qui est un passage à la catégorie des fractions, $\beta' $ le coréflecteur, et $\alpha$, $\beta$ les inclusions ; $\beta\, u_0 : E_0 \to F$ est pleinement fidèle (condition~4)).
\end{enonce}

\begin{proof}
\emph{Condition~1).} $Y = vP$ est dans $E_0$ si et seulement si $\eta_{vP}$ est inversible, et $uX \in F_0$ si et seulement si $\varepsilon_{uX}$ l'est : c'est l'énoncé sur les adjonctions idempotentes.

\emph{Le réflecteur.} Si $\operatorname{Im} u \subset F_0$, alors $vu(X) \in E_0$ pour tout $X$, par le second fait ; le foncteur $X \mapsto vu(X)$ arrive dans $E_0$. On prend pour unité $\eta_X : X \to vu(X)$ et pour coünité, en $Y \in E_0$, $\eta_Y^{-1} : vu(Y) \to Y$, naturelle parce que $\eta$ l'est. L'identité triangulaire du côté de $E_0$ est $\eta_Y\,;\eta_Y^{-1} = \mathrm{id}$. Celle du côté de $E$ demande $vu(\eta_X)\,;\eta_{vuX}^{-1} = \mathrm{id}$, c'est-à-dire $vu(\eta_X) = \eta_{vuX}$ ; c'est l'égalité établie dans la démonstration de (a)~$\Rightarrow$~(c), l'adjonction étant idempotente par cet énoncé. C'est le calcul de Hom de la page~5. Le coréflecteur se traite de même, avec la coünité $\varepsilon_P$, l'unité $\varepsilon_Q^{-1}$ en $Q \in F_0$, et l'égalité $uv(\varepsilon_P) = \varepsilon_{uvP}$ ; c'est le calcul entre crochets de la page~4.

\emph{Limites.} Un adjoint à gauche commute aux limites inductives, un adjoint à droite aux limites projectives. Une sous-catégorie réflexive a les limites inductives de la catégorie ambiante : on réfléchit la limite inductive calculée dans $E$.

\emph{Condition~4).} Pour $X \in E$, $u_0(vu(X)) = uvu(X)$, et $\varepsilon_{uX} : uvu(X) \to u(X)$ est inversible et naturel en $X$ : d'où $\beta\,u_0\,\alpha' \simeq u$. De même $\eta_{vP}^{-1} : vuv(P) \to v(P)$ donne $\alpha\,v_0\,\beta' \simeq v$. Le foncteur $\beta\,u_0$ est pleinement fidèle comme composé d'une équivalence et d'une inclusion pleine. Enfin, un adjoint à gauche dont l'adjoint à droite est pleinement fidèle est une localisation en les flèches qu'il inverse ; c'est le cas de $\alpha'$.
\end{proof}

\begin{enonce}[La descente d'adjoints, page 6]
Soient $\alpha : E_1 \to E$ pleinement fidèle d'adjoint à gauche $\alpha'$, et $\psi : E_1^{\circ} \to \mathrm{Ens}$ tel que $\psi \circ \alpha'$ soit représentable. Alors $\psi$ est représentable. Par suite, si $\beta : E_1 \to F$ est un foncteur tel que $u = \beta\alpha'$ admette un adjoint à droite, $\beta$ admet un adjoint à droite.
\end{enonce}

\begin{proof}
Soit $\xi \in E$ représentant $\psi \circ \alpha'$ : $\operatorname{Hom}_E(X, \xi) \simeq \psi(\alpha' X)$, naturellement en $X$. Si $f : X \to Y$ est dans $\Sigma$, la composition avec $f$, de $\operatorname{Hom}(Y, \xi)$ dans $\operatorname{Hom}(X, \xi)$, correspond à $\psi(\alpha' f)$, qui est bijective puisque $\alpha'(f)$ est inversible ; $\xi$ est donc $\Sigma$-local, et, par l'énoncé sur les objets locaux, il existe $\xi_1 \in E_1$ et un isomorphisme $\alpha \xi_1 \simeq \xi$. Pour $Z \in E_1$, on compose alors les bijections
\[
  \operatorname{Hom}_{E_1}(Z, \xi_1) \simeq \operatorname{Hom}_E(\alpha Z, \alpha \xi_1) \simeq \operatorname{Hom}_E(\alpha Z, \xi) \simeq \psi(\alpha'\alpha Z) \simeq \psi(Z),
\]
la première par la pleine fidélité de $\alpha$, la dernière par $\psi$ appliqué à la coünité $\alpha'\alpha Z \to Z$, inversible puisque $\alpha$ est pleinement fidèle. La naturalité en $Z$ résulte de celle de la coünité. Donc $\psi$ est représenté par $\xi_1$.

Pour la seconde assertion, $\beta$ admet un adjoint à droite si et seulement si, pour tout $P \in F$, le préfaisceau $Z \mapsto \operatorname{Hom}_F(\beta Z, P)$ sur $E_1$ est représentable. Composé avec $\alpha'$, c'est $X \mapsto \operatorname{Hom}_F(\beta\alpha' X, P)$, représentable puisque $\beta\alpha'$ a un adjoint à droite. La première assertion conclut.
\end{proof}

\begin{remarque}
L'exemple ci-dessus répond à la question de la marge de la page~3, avec $E'$ réduit à un objet ; la réponse de la lecture, \q{Elle l'est}, est juste.
\end{remarque}

\begin{remarque}
L'énoncé des pages~6 et~7 vaut pour toute sous-catégorie réflexive. La caractérisation de $\Sigma$ est classique, et l'identification des objets $\Sigma$-locaux à l'image essentielle, le \q{Or c'est vrai} de la page~7, est un lemme de base de la théorie des localisations.
\end{remarque}

\begin{remarque}
La cinquième condition de la page~3, \q{—— et $u'$ et $v'$ pl. fid.}, répète mot pour mot la deuxième ; l'énoncé ci-dessus en retient quatre.
\end{remarque}

\begin{remarque}
La descente de la page~6 tient, et sans l'hypothèse que $\beta$ soit pleinement fidèle : la démonstration ne s'en sert pas. Pour une catégorie $F$ quelconque, la réduction de la page au cas $\mathrm{Ob}\,F = \mathrm{Ob}\,E_1 \cup \{a\}$ revient à traiter un objet $P$ de $F$ à la fois, comme on l'a fait. La vérification suppose que les ensembles de flèches de $E$, $E_1$ et $F$ vivent dans le même univers, ce que demande la notion de foncteur représentable.
\end{remarque}

\begin{remarque}
La page~5 construit le réflecteur de $E_0$ sous la seule hypothèse $\operatorname{Im} u \subset F_0$ et le coréflecteur de $F_0$ sous $\operatorname{Im} v \subset E_0$ ; ces deux hypothèses sont équivalentes par l'énoncé sur les adjonctions idempotentes, et l'identité triangulaire du réflecteur passe par cette équivalence.
\end{remarque}

\noindent Ne sont pas démontrés ici : la reconstruction d'une adjonction idempotente à partir d'une sous-catégorie réflexive, d'une coréflexive et d'une équivalence entre elles, dans le sens qui part de ces données (page~5) ; aux pages~9 à~19, la propriété universelle de la catégorie monoïdale symétrique libre $\Phi(A)$ et le produit non parenthésé, les contractions et leur calcul graphique\piste{161-1-contractions-string-calculus}, le signe $\varepsilon(L)$ et les catégories de Picard, les théories de la page~13, l'extension de Kan et sa préservation des colimites (pages~17 et~19), le critère de Giraud et le calcul de densité.
